\documentclass[11pt]{article}

\usepackage[utf8]{inputenc}
\usepackage[T1]{fontenc}
\usepackage[margin=1in]{geometry}
\usepackage{amsmath,amssymb,amsthm}
\usepackage{booktabs}
\usepackage{microtype}
\usepackage{hyperref}
\usepackage[numbers,sort&compress]{natbib}
\usepackage{xcolor}

\hypersetup{
  colorlinks=true,
  linkcolor=blue!60!black,
  urlcolor=blue!50!black,
  citecolor=green!45!black,
}

\newtheorem{theorem}{Theorem}
\newtheorem{lemma}[theorem]{Lemma}
\newtheorem{corollary}[theorem]{Corollary}
\newtheorem{proposition}[theorem]{Proposition}
\theoremstyle{definition}
\newtheorem{definition}[theorem]{Definition}
\newtheorem{observation}[theorem]{Observation}
\theoremstyle{remark}
\newtheorem{remark}[theorem]{Remark}

% Monospace Lean identifiers, with a break opportunity after each "_" so long
% theorem names (e.g. insert_matching_common_order) never overflow a line.
\newcommand{\leanname}[1]{%
  \texttt{\begingroup\def\_{\textunderscore\allowbreak}#1\endgroup}}
\newcommand{\mcso}{\mathrm{MCS}}
\newcommand{\Adj}{\mathsf{adj}}
\newcommand{\neigh}{\mathsf{neigh}}
\setlength{\emergencystretch}{3em}

\title{Locality and stability of maximum cardinality search orderings
under edge updates, with a machine-checked implementation}

\author{Md Thoriqul Islam Thonmay \and Gregory Morse\\[6pt]
\normalsize Department of Programming Languages and Compilers,\\
Faculty of Informatics, E\"otv\"os Lor\'and University, Budapest.}

\date{October 2026}

\iffalse
ARXIV-METADATA-BEGIN
Title: Locality and stability of maximum cardinality search orderings under edge updates, with a machine-checked implementation
Authors: Md Thoriqul Islam Thonmay, Gregory Morse
Abstract: We study the problem of maintaining a maximum cardinality search (MCS) vertex ordering of an undirected graph under a single-edge update (insertion or deletion), instead of recomputing it from scratch. Every ordering we maintain is the canonical MCS ordering under smallest-index tie-breaking, and every update either leaves the ordering unchanged or re-greedily rebuilds one suffix. We give a machine-checked proof, in the Lean 4 theorem prover over Mathlib, of the update correctness, of the exact locality of a single-edge update (the first position whose choice can change lies in a window between the two endpoints; for a deletion, at a single position), of the refinement of a boolean-adjacency executable implementation to the abstract update, and of an operational unit-cost bound along the real control flow. The central structural result is a stability dichotomy: the set-valued MCS orderings of two graphs intersect whenever the two graphs track one another at distance at most one along every prefix; inserting or deleting a matching, or flipping a single edge, always preserves a common MCS ordering, while no common ordering survives inserting a triangle or a P_4, or a mixed matching update. All positive results and the three certified obstructions are proved in Lean and kernel-rechecked. To our knowledge this is the first maintenance of a maximum cardinality search ordering under edge updates, and the first machine-checked correctness proof of a maintained graph-search ordering.
Comments: 8 pages, 1 table. Code and proofs: https://github.com/thonmay/thesis-lean . Palomar registry: PALOMAR-2026-10-04-000009
Subjects: Data Structures and Algorithms (cs.DS); Combinatorics (math.CO)
MSC classes: 68Q25 (Primary) 05C85 (Secondary)
ACM classes: F.2.2; G.2.2; F.3.2
Report number: N/A -- leave blank
Journal reference: N/A -- leave blank
DOI: N/A -- leave blank
License: CC BY 4.0
Submitter is author: Yes
External artifacts: GitHub https://github.com/thonmay/thesis-lean; release/tag mcs-arxiv-v2; commit a92df358739288c32485b4039d4899b7fb04eaaa; Palomar source commit a872fddca7025ae0ecedd654238fa0c06729dc2e; verified 2026-10-04
AI disclosure: Acknowledgements and AI disclosure -- This work was developed with AI assistance in accordance with the ELTE IK policy on the use of AI tools. Three models, each driven through an agent harness over Lean~4, participated. DeepSeek-V4-Flash-0731 (provider DeepSeek; model identifier deepseek-ai/DeepSeek-V4-Flash-0731; agentic tool mode) carried the original development phase: the core definitions, the probe-then-repair updates, and their validity proofs, checked with lake build over Mathlib. Anthropic Claude Opus 5.5 (provider Anthropic; CLI coding-agent tool mode) planned and drafted the resubmission, wrote the single-source canonical definitions, the locality and ExecUpdate modules, the stability module, and the unit-cost instrumentation, and ran the brute-force stability experiments. Qwen3.8-27B-FP8 (provider Alibaba; model identifier Qwen/Qwen3.8-27B-FP8; agentic tool mode) repaired the executable refinement and unit-cost proof bodies, contributed the certified stability counterexamples, and wired the comparator. In every phase the model's proof state was machine-checked with lake build and independently re-checked by the leanchecker external kernel; the transport and bridge lemmas were written by hand. The human authors reviewed the resulting definitions, statements, proofs, experiments, and manuscript text and assume full responsibility for the final work. No AI system is an author.
ARXIV-METADATA-END
\fi
\begin{document}
\maketitle

\begin{abstract}
We study the problem of maintaining a \emph{maximum cardinality search} (MCS)
vertex ordering of an undirected graph under a single-edge update (insertion or
deletion), instead of recomputing it from scratch. Every ordering we maintain is
the canonical MCS ordering under smallest-index tie-breaking, and every update
either leaves the ordering unchanged or re-greedily rebuilds one suffix. We give
a machine-checked proof, in the Lean~4 theorem prover over Mathlib, of the update
correctness, of the exact locality of a single-edge update (the first position
whose choice can change lies in a window between the two endpoints; for a
deletion, at a single position), of the refinement of a boolean-adjacency
executable implementation to the abstract update, and of an operational unit-cost
bound along the real control flow. The central structural result is a
\emph{stability dichotomy}: the set-valued MCS orderings of two graphs intersect
whenever the two graphs track one another at distance at most one along every
prefix; inserting or deleting a \emph{matching}, or flipping a single edge, always
preserves a common MCS ordering, while no common ordering survives inserting a
triangle or a $P_4$, or a mixed matching update. All positive results and the
three certified obstructions are proved in Lean and kernel-rechecked. To our
knowledge this is the first maintenance of a maximum cardinality search
\emph{ordering} under edge updates, and the first machine-checked correctness
proof of a maintained graph-search ordering.
\end{abstract}

\section{Introduction}
\label{sec:intro}

A \emph{maximum cardinality search} (MCS) ordering, introduced by Tarjan and
Yannakakis~\cite{tarjan1984mcs}, numbers the vertices of a graph so that each
vertex is chosen to maximize the number of already-chosen neighbours. Reverse MCS
orderings are the canonical tool for chordal-graph recognition and give perfect
elimination orderings~\cite{rose1976algorithmic,tarjan1984mcs}; MCS is a simple
restriction of the lexicographic breadth-first search (LexBFS) of Rose, Tarjan and
Lueker~\cite{rose1976algorithmic}, and the two ordering families are incomparable
even on chordal graphs~\cite{corneil2008unified}. This paper asks a dynamic
question about MCS orderings: given an MCS ordering of a graph, how much of it
must change after a single edge is inserted or deleted?

Our answer is structural rather than asymptotic. We make \textbf{no asymptotic
separation claim}: in the unit-cost model we adopt, the worst-case update cost is
at most $n$ (insertion) and $n+1$ (deletion) units, the same order as recomputing
from scratch, and the static MCS of Tarjan--Yannakakis is linear time. Instead we
prove three things. First, a \emph{stability dichotomy} (Section~\ref{sec:stab})
that identifies exactly where the two graphs stop sharing an MCS ordering.
Second, the \emph{locality} of a single-edge update (Section~\ref{sec:locality}):
the window, and for a deletion the single point, in which an update can first
disturb the ordering. Third, a \emph{verified executable} (Section~\ref{sec:exec})
and an \emph{operational unit cost} (Section~\ref{sec:cost}). Every theorem is
proved in Lean~4 and independently re-checked by an external kernel checker;
Section~\ref{sec:formalization} describes the artifact.

\subsection{Contributions}

\begin{enumerate}
\item \textbf{Stability dichotomy.} For which edge-modification sets $F$ is
$\mcso(G)\cap\mcso(G\triangle F)\neq\varnothing$, where $\mcso(\cdot)$ denotes
the \emph{set} of MCS orderings under arbitrary tie-breaking? We prove the
intersection is nonempty for every insertion-only and deletion-only matching, and
that the single-edge flip always preserves a common ordering. We certify that the
two hypotheses of the positive theorem cannot be relaxed: a triangle and a $P_4$
produce no common ordering (the ``$+1$'' bound fails), and a mixed matching
update produces none (the one-directional condition is necessary). An online
corollary (\leanname{common\_order\_extends\_prefix}) shows any common MCS prefix
extends to a full common ordering without backtracking, tying the set-valued theory
to the constructive update.
\item \textbf{Locality.} A single-edge insertion can first change the ordering
only inside the positional window between the two endpoints; a deletion only at
the later endpoint's position. When no break fires the ordering is \emph{literally
unchanged}, not merely still valid.
\item \textbf{Verified executable and unit cost.} A boolean-adjacency
implementation refines the abstract update; an instrumented version counts
legality checks and pick steps along the real control flow, with an exact bound
of $1$ unit for a no-break deletion and at most the window size for a no-break
insertion.
\end{enumerate}

\paragraph{Honesty about the cost model.} One cost unit equals one legality
check or one pick step, and a unit is itself $\Theta(n\cdot|\text{chosen}|)$ work
in this implementation (a pick filters the whole chosen list per candidate). The
unit-cost model therefore compares our update against from-scratch recomputation
\emph{within the same model}, and says nothing about wall-clock parity with the
linear-time static algorithm of Tarjan and Yannakakis~\cite{tarjan1984mcs}.

\section{Preliminaries}
\label{sec:prelim}

Fix $n\in\mathbb N$ and write $V_n=\mathrm{Fin}\,n$ for the vertex set. A graph
is a symmetric, loop-free binary relation $\Adj\colon V_n\times V_n\to\mathsf{Prop}$
equipped with proofs $\Adj\,u\,v\to\Adj\,v\,u$ and $\lnot\Adj\,v\,v$. For a set
$S\subseteq V_n$ and a vertex $w$, write $\neigh_G(S,w)$ for the number of
neighbours of $w$ inside $S$. A vertex $w\notin S$ is a \emph{legal MCS pick} for
the chosen set $S$ if $\neigh_G(S,w)$ is maximum over $V_n\setminus S$; it is a
\emph{canonical} pick if it is additionally smallest-index among the maxima. An
ordering $v_0,\dots,v_{n-1}$ is an MCS ordering if every prefix picks legally.

\begin{definition}[Canonical vs.\ set-valued]
\leanname{IsMCSOrdering} is the canonical predicate: at each step the picked
vertex is an $\mcso$-maximum with the smallest index. \leanname{IsMCSOrderingAny}
drops the smallest-index clause and asserts only that each step is
\emph{some} $\mcso$-maximum. The canonical ordering of a graph is unique
(\leanname{IsMCSOrdering\_unique}); $\mcso(G)$ denotes the \emph{set} of
orderings satisfying \leanname{IsMCSOrderingAny}.
\end{definition}

A single-edge update is modelled by $\leanname{FlipOf}\ G\ G'\ u\ v$: $G'$ is $G$
with the edge $\{u,v\}$ toggled (inserted if absent, deleted if present). The
executable layer works over a concrete boolean adjacency matrix
$\leanname{Adj}\ n = \mathsf{List}\ (\mathsf{List}\ \mathsf{Bool})$ with a
transport \leanname{toUGraph}.

\section{The stability dichotomy}
\label{sec:stab}

The main structural question is: when do $G$ and $G'$ admit a \emph{common} MCS
ordering? The positive answer is stated under arbitrary tie-breaking, so it is
about the sets $\mcso(G)$ and $\mcso(G')$.

\begin{theorem}[Common ordering, \leanname{exists\_common\_ordering}]
\label{thm:common}
Let $G,G'$ be graphs on $n$ vertices such that for every $S\subseteq V_n$ and
every $x\in V_n$,
\[
\neigh_G(S,x)\ \le\ \neigh_{G'}(S,x)\ \le\ \neigh_G(S,x)+1.
\]
Then $\mcso(G)\cap\mcso(G')\neq\varnothing$.
\end{theorem}

The proof is an induction that builds the common ordering one vertex at a time.
Its key step is that, under the hypothesis, every chosen set $S\neq V_n$ admits a
vertex that is an MCS maximum in \emph{both} graphs simultaneously
(\leanname{exists\_common\_pick}).

\begin{corollary}[Online form, \leanname{common\_order\_extends\_prefix}]
\label{cor:online}
The induction in Theorem~\ref{thm:common} is stronger than the theorem: any
common MCS prefix---a duplicate-free partial ordering whose every step is a legal
MCS step in \emph{both} graphs---extends to a full common MCS ordering of both
graphs. A maintainer that, at each step, picks \emph{any} vertex that is maximum in
both graphs can therefore extend its prefix forever, with no lookahead and no
backtracking. This is a separate entry in the verified artifact, proved by strong
induction on the number of unplaced vertices, and it is the statement that makes the
canonical update of Section~\ref{sec:exec} one realization of the constructive
half.
\end{corollary}

We apply Theorem~\ref{thm:common} to the update model. A \emph{matching} is a
graph in which every vertex has degree at most one.

\begin{corollary}[Matching insertions/deletions, \leanname{insert\_matching\_common\_order}, \leanname{delete\_matching\_common\_order}]
\label{cor:matching}
If $G'$ is obtained from $G$ by inserting a matching, or by deleting a matching,
then $\mcso(G)\cap\mcso(G')\neq\varnothing$.
\end{corollary}

\begin{corollary}[Single-edge flip, \leanname{flip\_common\_order}]
A single-edge flip always preserves a common MCS ordering:
for any $u\neq v$ with $\leanname{FlipOf}\ G\ G'\ u\ v$, $\mcso(G)\cap\mcso(G')\neq\varnothing$.
\end{corollary}

Both hypotheses of Theorem~\ref{thm:common} are precisely what one would try to
relax, and both fail. The ``$+1$'' bound cannot become ``$+2$'', and the
assumption that the update is one-directional (all inserts or all deletes) cannot
be dropped. Each negative result is a concrete small instance, proved by a pruned
search whose soundness is itself a Lean theorem (\leanname{no\_common\_order\_of\_search}
and \leanname{commonSearch\_of\_common}) and discharged by the kernel on an
explicit edge list (not by \leanname{native\_decide}).

\begin{theorem}[Certified obstructions]
\label{thm:obstructions}
Certified by kernel \leanname{decide} on explicit edge lists (not
\leanname{native\_decide}):

\begin{itemize}
\setlength\itemsep{2pt}
\item \textbf{Triangle} (\leanname{triangle\_no\_common\_order}). On $n=5$,
inserting a triangle into $G=\{(0,2),(0,4),(1,2),(1,3)\}$ (adding
$(2,3),(2,4),(3,4)$) leaves no common MCS ordering:
$\mcso(G)\cap\mcso(G')=\varnothing$. The triangle has maximum $F$-degree $2$.

\item \textbf{Path} (\leanname{p4\_no\_common\_order}). On $n=6$, inserting the
triangle-free path $2\text{--}5\text{--}4\text{--}3$ into
$G=\{(0,3),(0,5),(1,2),(1,4),(2,3)\}$ also leaves no common ordering. Unlike the
triangle, the inserted $P_4$ is triangle-free, so the obstruction does not reduce
to the presence of a clique; both cases have maximum $F$-degree $2$ and witness
the same ``inflate by two'' failure, on a clique and on a path respectively.

\item \textbf{Mixed matching} (\leanname{mixed\_matching\_no\_common\_order}). On
$n=4$, flipping the matching $\{01,23\}$ as one insertion and one deletion
(deleting $01$, inserting $23$) into $G=\{(0,1),(0,3),(1,2)\}$ leaves no common
ordering. This isolates the one-directional condition as necessary.
\end{itemize}
\end{theorem}

The positive theorem is \emph{sufficient, not necessary} (see
Section~\ref{sec:conclusion}): a single insertion can inflate one vertex's count
by two and still leave a common ordering, so a violation of the ``$+1$'' clause
does not automatically destroy commonness. The exact boundary is posed as an open
problem in Section~\ref{sec:conclusion}.

\section{Locality of a single-edge update}
\label{sec:locality}

The canonical update is the \emph{probe-then-repair} procedure: scan the current
ordering for the first step whose old choice stops being a legal MCS pick in the
updated graph, and re-greedily rebuild only the suffix from that point. Its
correctness is the local counterpart of the dichotomy.

\begin{theorem}[Update correctness, \leanname{init\_valid}, \leanname{insert\_update\_valid}, \leanname{delete\_update\_valid}]
Starting from a valid canonical ordering, a single insertion or deletion returns
a valid canonical MCS ordering of the updated graph; the empty-prefix regreedy
initializes correctly (\leanname{greedySuffix\_full}) and updates preserve the
legal prefix (\leanname{update\_from\_prefix}).
\end{theorem}

\begin{theorem}[Update equals recompute, \leanname{insert\_update\_eq\_initOrder}, \leanname{delete\_update\_eq\_initOrder}]
The returned ordering is exactly the from-scratch canonical ordering of the
updated graph; hence (with \leanname{IsMCSOrdering\_unique}) the maintained state
is the canonical ordering, not just some valid one.
\end{theorem}

The locality of the probe is where the algorithm gains over recomputation.

\begin{theorem}[Insertion window, \leanname{legal\_upto\_earlierPos}, \leanname{legal\_after\_laterPos}, \leanname{insert\_break\_iff}]
For an insertion of $\{u,v\}$ with $u$ earlier than $v$ in the canonical order,
every position up to $u$'s position and every position after $v$'s position keeps
its old choice legal; inside the window, a step breaks exactly when the later
endpoint $v$ now beats the old choice at that step.
\end{theorem}

\begin{theorem}[Deletion single-point, \leanname{delete\_legal\_between}]
For a deletion of $\{u,v\}$, only the later endpoint's own position can be the
first break; positions strictly between the endpoints stay legal.
\end{theorem}

\begin{theorem}[Null update, \leanname{insert\_update\_eq\_self\_iff}, \leanname{delete\_update\_eq\_self\_iff}]
The insertion update, respectively the deletion update, is the identity exactly
when the ordering is already the canonical ordering of the updated
graph---otherwise it regreedies a suffix.
\end{theorem}

\section{Verified executable}
\label{sec:exec}

To rule out the gap where a ``specification'' is proved but nothing executable
implements it, the development carries a concrete boolean-adjacency
implementation and proves it refines the abstract update.

\begin{theorem}[Executable refinement, \leanname{exec\_insert\_update\_eq}, \leanname{exec\_delete\_update\_eq}]
The executable insertion and deletion functions on $\leanname{Adj}\ n$ return, up to
the vertex-value map, exactly the abstract update. The flip of the matrix preserves
symmetry and loop-freeness (\leanname{setEdge\_preserves\_symm},
\leanname{setEdge\_preserves\_loopless}), and the executable from-scratch MCS
equals the abstract greedy suffix (\leanname{mcsOrder\_eq\_greedySuffix}).
\end{theorem}

The three theorems \leanname{exec\_insert\_recompute\_valid},
\leanname{exec\_delete\_recompute\_valid}, and \leanname{mcsOrder\_eq\_greedySuffix}
deliberately certify \emph{recomputation}, not the dynamic update: they state that
recomputing the MCS of the flipped matrix from scratch is valid. The dynamic
update is certified by the refinement theorems above, which take the honest
names.

\section{Unit cost along the real control flow}
\label{sec:cost}

An instrumented variant of the executable update counts, at each step, one unit
for a legality check and one unit for a pick step, and returns the count together
with the ordering. Its first projection is provably the plain update
(\leanname{exec\_insert\_update\_c\_fst}, \leanname{exec\_delete\_update\_c\_fst}),
so the bounds below attach the cost to the update that is actually run.

\begin{theorem}[Sharp bounds]
\label{thm:cost}
Let $n=|V|$. In the unit-cost model:

\begin{enumerate}
\setlength\itemsep{2pt}
\item \textbf{No-break deletion} (\leanname{exec\_delete\_cost\_of\_no\_break}):
if a deletion's probe finds no break, the cost is \emph{exactly} $1$ unit.

\item \textbf{No-break insertion} (\leanname{exec\_insert\_cost\_of\_no\_break}):
if an insertion's probe finds no break, the cost is at most the window size
$k=\mathrm{pos}(v)-\mathrm{pos}(u)$ for the later endpoint $v$ (and
$\mathrm{pos}(u)-\mathrm{pos}(v)$ if $u$ is later).

\item \textbf{Break cost} (\leanname{exec\_insert\_cost\_of\_break},
\leanname{exec\_delete\_cost\_of\_break}): a break at position $k$ costs the probe
plus at most $n-k$ regreedy steps.

\item \textbf{Worst case} (\leanname{exec\_insert\_cost\_le},
\leanname{exec\_delete\_cost\_le}): at most $n$ (insertion) and $n+1$ (deletion)
units---matching, within this model, from-scratch recomputation.
\end{enumerate}
\end{theorem}

The two genuinely sharp statements are the first two: a no-break deletion costs
one check, a no-break insertion costs a probe of the window. These are the
operational content of the locality results of Section~\ref{sec:locality}.

\section{The Lean formalization}
\label{sec:formalization}

The development is a Lean~4 / Mathlib project (toolchain pinned to
4.35.0-rc2), organized as a Palomar registry challenge and
\emph{registered} as Palomar entry \leanname{PALOMAR-2026-10-04-000009}
(version~1, status \texttt{registered}).\footnote{Entry:
\url{https://palomar-registry.org/entry.html?id=PALOMAR-2026-10-04-000009&version=1};
machine-readable record:
\url{https://data.palomar-registry.org/entries/PALOMAR-2026-10-04-000009-v1.json}.}
A module
\leanname{Challenge} poses the headline statements (with explicit \texttt{sorry}
bodies and canonical definitions in Mathlib-only code), and a module
\leanname{Solution} proves them; a single source of truth generates both so
every posed statement is textually identical to its proved counterpart. The
challenge/solution comparator lists $37$ theorem names, five to eight per category,
all of which are introduced inline in Sections~\ref{sec:stab}--\ref{sec:cost}:

\begin{center}
\begin{tabular}{lrl}
\toprule
Category & count & Section \\
\midrule
correctness & $5$ & \ref{sec:locality} \\
uniqueness / update = recompute & $3$ & \ref{sec:locality} \\
locality & $6$ & \ref{sec:locality} \\
executable bridge & $5$ & \ref{sec:exec} \\
executable refinement & $2$ & \ref{sec:exec} \\
unit cost & $8$ & \ref{sec:cost} \\
stability (incl.\ online corollary and $3$ obstructions) & $8$ & \ref{sec:stab} \\
\midrule
total & $37$ & \\
\bottomrule
\end{tabular}
\end{center}

Every theorem audits to only the three standard Lean axioms (\texttt{propext},
\texttt{Quot.sound}, \texttt{Classical.choice}): there is no \texttt{sorry} outside
\leanname{Challenge}, and no \leanname{native\_decide}. An \texttt{AxiomAudit} module
prints the axioms of each comparator theorem, enforced as a CI step; a
definitions-sync check enforces byte-identity of the canonical definitions; and an
external kernel checker (\texttt{leanchecker}) re-validates each module. The
certified obstructions are independently cross-checked by a separate exhaustive
Python enumerator that reproduces each witness edge list exactly.

\paragraph{Artifact and reproducibility.} The complete Lean development, its
comparator, the axiom audit, and the empirical verification scripts are public
at \url{https://github.com/thonmay/thesis-lean}, pinned to the annotated tag
\leanname{mcs-arxiv-v2} (commit \texttt{a92df35}). The Palomar registration
record was computed from the Lean proof at commit \texttt{a872fdd}, which is an
ancestor of the tagged artifact. The repository README maps
every theorem named in this paper to its Lean source file and gives the exact
build and verification commands (\texttt{lake build}, the \texttt{AxiomAudit}
step, and per-module \texttt{leanchecker}). The formalization is
\emph{logically required} for Sections~\ref{sec:stab}--\ref{sec:cost}: every
positive statement and every certified obstruction is a theorem in this
repository. Three Python scripts are \emph{supporting} (not proof premises):
the exhaustive stability enumerator
(\leanname{scripts/mcs\_stability.py}) reproduces the counterexample witness edge
lists; \leanname{scripts/incomparability\_witness.py} enumerates both MCS and
LexBFS ordering families over all tie-breaks on the witness of
Section~\ref{sec:related} and certifies the $204/204/40$ count; and
\leanname{scripts/mcs\_highn\_crosscheck.py} runs the adopted candidate
implementation (\leanname{scripts/candidates/gen003\_c02.py}) on random graphs up
to $n=200$, checking after every edge flip that its ordering is a legal MCS
ordering and equals a from-scratch recomputation.

\section{Related work}
\label{sec:related}

The nearest dynamic prior art maintains a property of a \emph{chordal} graph via
an MCS-adjacent object, rather than the MCS ordering of an arbitrary graph
itself. Ibarra~\cite{ibarra2008dynamic} maintains a clique-tree representation of
a chordal graph under insertions and deletions; Banerjee--Raman--Satti
\cite{banerjee2018maintaining} maintain a perfect elimination ordering under
deletions only, with a different window notion and a cell-probe lower bound for
dynamic chordality; Mezzini~\cite{mezzini2012fully} gives $O(1)$-query/$O(n^2)$-update
chordal maintenance. None of these maintains the MCS ordering itself, and none is
machine-checked. Outside MCS, a maintained BFS ordering under edge updates appears
in Morse--Kozsik~\cite{morse2026bfs} (a framework for fully dynamic BFS spanning
trees and their induced ordering in directed graphs), with a separate
dynamic-LexBFS implementation by the same author; our result is specific to MCS
orderings of undirected graphs, and is the only one machine-checked. For verified
dynamic graph algorithms, the closest is Gu\'eneau et al.~\cite{gueneau2019formal}
(incremental cycle detection in Coq with time credits); no prior formalization of
MCS or of a perfect elimination ordering, and no machine-checked fully-dynamic
search-ordering maintenance, was found in any proof assistant. MCS-M
\cite{berry2004mcs} is a multi-sweep extension of MCS for computing minimal
triangulations; it is static and does not maintain an ordering under updates. The
incomparability of MCS and LexBFS orderings is due to
Corneil--Krueger~\cite{corneil2008unified}; a brute-force enumeration of the
incomparability witness (a chordal 7-vertex graph, all tie-breaks) checks out to
the digit: $204$ MCS orderings, $204$ LexBFS orderings, a symmetric difference of
$40$ on each side. Brandst\"adt, Dragan and Nicolai~\cite{brandstadt1997lexbfs}
show a common-perfect-elimination-ordering result for LexBFS powers that fails for
MCS, which is adjacent but negative prior art for our positive stability statement.

\section{Conclusion and open problems}
\label{sec:conclusion}

We have given a machine-checked account of maintaining the canonical MCS ordering
under single-edge updates: exact locality, an executable refinement, a unit-cost
bound, and a sharply delimited stability dichotomy for the set-valued MCS
orderings. Beyond registering the formalization, three directions stand out.

\paragraph{The stability boundary.} Conjecture (Star--Matching): for a one-sided
update $F$, $\mcso(G)\cap\mcso(G\triangle F)\neq\varnothing$ for every host graph
$G$ iff $F$ is a star plus a matching, i.e. at most one vertex has $F$-degree
$\ge 2$. The matching side is proved here (Corollary~\ref{cor:matching}); standard
combinatorial reasoning shows a one-sided update $F$ with two vertices of
$F$-degree $\ge 2$ always contains a triangle, a $P_4$ or a $2P_3$ as a subgraph,
of which the first two are certified here (Theorem~\ref{thm:obstructions}) and the
third is not; the star case needs a non-local inflation argument and is open.

\paragraph{Sharpenings of the probe.} Two experimentally verified refinements
remain to be formalized: (i) an insertion never breaks at the later endpoint's own
position, so the window is genuinely half-open $(u,v]$ and the probe scans one
position too many today; and (ii) the break criterion can be stated against the
\emph{old} graph (a step breaks iff the later endpoint was already tied, or one
below with a smaller index), making the probe decidable from the old counts. Both
were observed exhaustively and would tighten the locality statements.

\paragraph{Sublinear probing.} A separate (not machine-checked) investigation
suggests the first-break predicate is not a prefix over the window, which would
rule out a segment-tree/Fenwick binary-search probe for exact MCS-ordering
maintenance and leave the linear probe the right shape. Documenting this
information-theoretic obstacle is a modest negative contribution alongside the
positive results.

\section*{Acknowledgements and AI disclosure}

This work was developed with AI assistance in accordance with the ELTE IK policy
on the use of AI tools. Three models, each driven through an agent harness over
Lean~4, participated. DeepSeek-V4-Flash-0731 (provider DeepSeek; model identifier
deepseek-ai/DeepSeek-V4-Flash-0731; agentic tool mode) carried the original
development phase: the core definitions, the probe-then-repair updates, and their
validity proofs, checked with lake build over Mathlib. Anthropic Claude Opus 5.5
(provider Anthropic; CLI coding-agent tool mode) planned and drafted the
resubmission, wrote the single-source canonical definitions, the locality and
ExecUpdate modules, the stability module, and the unit-cost instrumentation, and
ran the brute-force stability experiments. Qwen3.8-27B-FP8 (provider Alibaba;
model identifier Qwen/Qwen3.8-27B-FP8; agentic tool mode) repaired the executable
refinement and unit-cost proof bodies, contributed the certified stability
counterexamples, and wired the comparator. In every phase the model's proof state
was machine-checked with lake build and independently re-checked by the leanchecker
external kernel; the transport and bridge lemmas were written by hand. The human
authors reviewed the resulting definitions, statements, proofs, experiments, and
manuscript text and assume full responsibility for the final work. No AI system is
an author.

\bibliographystyle{plainnat}
\bibliography{references}

\end{document}